\documentclass[t,aspectratio=169]{beamer}

\usepackage{soul}
\usepackage{tikz}
\usepackage{multicol}
\usepackage[utf8]{inputenc}

% ----------------------------------------------------------------------------------------------------------------------------------------

\title{\includegraphics[width=45pt]{oepecsetlogo.png}\\(KTXFI1EBNF) Physics I. Practice}

\author{Dr.~Gábor FACSKÓ, PhD\\\tiny Assistant Professor / Senior Research Fellow\\\textit{facsko.gabor@uni-obuda.hu}}

\institute{\Tiny Óbuda University, Faculty of Electrical Engineering, 1084 Budapest, Tavaszmező u. 17.\\Wigner Research Centre for Physics, Department of Space Physics and Space Technology, 1121 Budapest, Konkoly-Thege Miklós út 29–33.\\ \url{https://wigner.hu/~facsko.gabor}}

\date{February 16,~2026}

\logo{\includegraphics[height=0.9 true cm]{kekcsik.png}}

% ----------------------------------------------------------------------------------------------------------------------------------------

\begin{document}

\frame{\titlepage}

% ----------------------------------------------------------------------------------------------------------------------------------------

\begin{frame}%[fragile,squeeze,allowframebreaks]
\frametitle{Exercises}
\begin{itemize}
\setlength{\parskip}{0pt}
\setlength{\itemsep}{0pt}
\item \textbf{Mechanics I:} Newton’s laws, equations of motion (review). Momentum. Work-energy theorem. Forms of mechanical energy.
\item \textbf{Mechanics II:} Weight, gravitational force, gravitational field. Work done by gravity. Conservative and dissipative force fields. Power.
\item Kinematics recap:
  \begin{itemize}
  \item Definitions: Path length, displacement, velocity, average velocity, acceleration.
  \item Uniform rectilinear motion: $\mathbf{r}(t) = \mathbf{r}_{0} + \mathbf{v} \cdot t$, $\mathbf{v} = \text{const}$, $\mathbf{a} = 0$.
  \item Uniformly accelerated rectilinear motion: $\mathbf{r}(t) = \mathbf{r}_{0} + \mathbf{v}_{0} \cdot t + \frac{1}{2}\mathbf{a}t^{2}$, $\mathbf{v}(t) = \mathbf{v}_{0} + \mathbf{a} \cdot t$, $\mathbf{a} = \text{const}$.
  \item Uniform circular motion: $\phi(t) = \omega \cdot t$, $\omega = \text{const}$, $\beta = 0$.
  \item Uniformly accelerated circular motion: $\phi(t) = \phi_{0} + \omega_{0} \cdot t + \frac{1}{2}\beta t^{2}$, $\omega(t) = \omega_{0} + \beta \cdot t$, $\beta = \text{const}$.
  \end{itemize}
\item Example: Projectile motion: \textit{At what angle will a projectile have the maximum range for a given initial speed $v$?}
\item \st{Non-conservative forces}
\end{itemize}
\end{frame}


% ----------------------------------------------------------------------------------------------------------------------------------------

\begin{frame}
%\frametitle{Minden jó, ha a vége jó}

\vfill
\begin{center}
\textbf{\huge The End}

\bigskip
Thank you for your attention!
\end{center}
\vfill
\textbf{Acknowledgements} Thank you for the course materials for Szilárd~ZSÓKA and Zoltán~VARGA.
\end{frame}

\end{document}
